% Part: applied-modal-logics
% Chapter: temporal-logic
% Section: properties-accessibility

\documentclass[../../../include/open-logic-section]{subfiles}

\begin{document}

\olfileid[pt]{aml}{tl}{acc}

\olsection{Propriedades dos frames temporais}

Como nossos modelos temporais não impõem nenhuma condição à
relação~$\prec$, o único de nossos axiomas conhecidos que vale em
todos os modelos é~$K$, ou seus análogos $K_{\Gtemp}$ e~$K_{\Htemp}$:
\begin{align*}
\tag{$K_{\Gtemp}$}  \Gtemp (p \to q) & \to (\Gtemp p \to \Gtemp q)\\
\tag{$K_{\Htemp}$}  \Htemp (p \to q) & \to (\Htemp p \to \Htemp q)
\end{align*}

Entretanto, se quisermos que nossos modelos imponham condições mais
rigorosas à representação do tempo, por exemplo garantindo que $\prec$
seja uma ordem linear, obteremos outras validades em nossos modelos.

\begin{table}[t]
  \begin{tabular}{| p{.48\textwidth} || p{.48\textwidth} |}
    \hline
    {\emph{Se $\prec$ é \dots}} & {\emph{então \dots é verdadeira em~$\mModel{M}$:}} \\
    \hline \hline
    \emph{transitiva}: \newline
    $\forall u \forall v \forall w ((u \prec v \land v \prec w) \lif u \prec w)$ & 
    $\Ftemp \Ftemp p \lif \Ftemp p$  \\
    \hline 
    \emph{linear}: \newline
    $\forall w \forall v (w \prec v \lor w = v \lor v \prec w)$ &  
    $(\Ftemp \Ptemp  p \lor \Ptemp \Ftemp p) \lif (\Ptemp p \lor p \lor \Ftemp p)$\\
    \hline
    \emph{densa}: \newline
    $\forall w \forall v (w \prec v \to \exists u(w \prec u \land u \prec v))$ &  
    $\Ftemp p \lif \Ftemp \Ftemp p$ \\
    \hline
    \emph{ilimitada (passado)}: \newline
    $\forall w \exists v( v \prec w)$ &  
    $\Htemp p \to \Ptemp p$ \\
    \hline
    \emph{ilimitada (futuro)}: \newline
    $\forall w \exists v( w \prec v)$ &  
    $\Gtemp p \to \Ftemp p$ \\
    \hline
  \end{tabular}
  \caption{Algumas propriedades de correspondência de frames temporais.}
  \ollabel{tab:correspondence}
\end{table} 

Várias propriedades de \olref{tab:correspondence} podem parecer
características desejáveis para um modelo destinado a representar o
tempo. Contudo, vale observar que, embora possamos impor à relação
$\prec$ quaisquer condições que quisermos, nem todas elas correspondem
a !!{formula}s que possam ser expressas na linguagem da lógica
temporal. A irreflexividade, por exemplo, isto é, a ideia de que
$\forall w \lnot (w \prec w)$, não possui uma fórmula correspondente
na lógica temporal. 

\end{document}
